% Section 5: Experiments
% Source: ../../drafts/Paper_论文向总览.md §8 (experiment design) +
% doc/ClawEval_Metadata.md (benchmark) + doc/CL_Update_Sunhao.md (20-experiment route).

\section{Experiments}
\label{sec:exp}

\subsection{Setup}

\paragraph{Benchmark and metric.}
We evaluate on \textsc{ClawEval}, an agentic benchmark of $300$ tasks over three
splits (General $161$ / Multimodal $101$ / Multi-turn $38$); we use the
text-only subset of $195$ tasks. Each task is scored as
\begin{equation}
  \text{score} = s_{\text{safety}}\times\big(0.8\cdot s_{\text{completion}} + 0.2\cdot s_{\text{robustness}}\big),
\end{equation}
under a \textbf{Pass\textsuperscript{3}} criterion (all three independent runs must
pass). Crucially, the training reward and the evaluation metric are the same
model-judge scoring function, so reward and evaluation are structurally
consistent. The 27B policy is deployed on 64 GPUs under a fully-asynchronous
40-inference / 24-training split.

\paragraph{Ablation route.}
We run a suite of controlled ablations organized as a dependency tree, each
comparison changing exactly one variable so that a component's marginal
contribution is attributable:
\begin{itemize}
  \item \textbf{Phase 1 (B1)}: pure RL, establishing the forgetting lower bound
    (no replay, no KL; entropy always on).
  \item \textbf{Phase 2 (KL)}: $\lambda_2$ sweep and the $\pi_{ref}$ anchor choice
    ($\pi_0$ vs.\ $\pi_{t-1}$).
  \item \textbf{Phase 3 (Replay)}: bucket structure vs.\ CLEAR~\cite{A2}; priority
    type (anti-forgetting vs.\ reward); U-shaped block weighting (W0 with
    $\gamma=\delta=1$ vs.\ W2 with $\gamma=\delta=0.88$); buffer capacity.
  \item \textbf{Phase 4 (KL $\times$ Replay)}: the $2\times2$ combination of the
    two consolidation mechanisms.
  \item \textbf{Phase 5 (rollout scale)}: $1024\times8 \to 4096\times8$.
\end{itemize}
Each pairwise contrast isolates a single factor (e.g., R3$\to$R4 toggles priority,
R4$\to$R5 swaps anti-forgetting for reward priority, R4$\to$R4-w adds the U-shaped
weight). The multi-turn query-construction pipeline is not given a standalone
ablation---its effect is validated through process monitoring and the ClawEval
Multi-turn split as an endpoint metric.

\paragraph{Process metrics (Echo-Trap / forgetting early warning).}
Alongside endpoint scores we log per-step diagnostics: output-entropy curves (a
$>\!50\%$ drop within the first 100 steps triggers a larger $\lambda_4$),
trajectory diversity (within-group distinct-$n$ / self-BLEU), the KL trend, the
$L_{replay}/L_{rl}$ ratio, per-loss gradient norms, and per-bucket buffer dynamics
(written to a sidecar by \texttt{trainer/replay\_metrics.py}).

\subsection{Results}

\TODO{headline forgetting / CL-Score numbers vs.\ CLEAR and pure-RL baselines;
per-phase ablation deltas; process-metric curves. Pending 64-GPU cluster runs.}

